\section{Conformance and reproducibility}\label{sec:conformance}
This section defines the evidence needed for implementation review. A conforming evaluator \must{} identify the profile and supported domain and agree with its objective and feasibility semantics. It \must{} fail explicitly on unsupported identities, malformed instances, missing dependencies and arithmetic/resource errors. The order or wording of reported validation errors is not standardized here. An invalid or inconclusive result carries no accepted authoritative objective.

The static companion tests cover original-token parsing, binary64 versus decimal interpretation, integer/grid eligibility, structural coverage, route-list ordering, flat versus grouped rounding, depot/service/window boundaries, overflow, exact result encoding and display collisions. They include exact midpoint quantization and a large-value witness where sequential binary64 loses increments. These synthetic fixtures deliberately exercise boundaries that observed incumbents may not reach. They do not replace a complete family audit or certify a native implementation's full input domain.

\begin{table}[htbp]
\small
\begin{tabularx}{\linewidth}{@{}>{\raggedright\arraybackslash}p{5cm}XX@{}}
\toprule
Synthetic operands or condition & Exact evaluation & Other declared evaluation\\
\midrule
Binary64 $2^{53},1,1$ in one flat sum & $2^{53}+2$ & Sequential flat returns $2^{53}$.\\
Same travel sequence, final deadline $2^{53}$ & Exact-binary infeasible & Rounded forward timing accepts.\\
Tokens \id{0.33}, \id{0.56}, \id{0.11}, final deadline 1 & Exact-binary infeasible & Exact-decimal accepts at 1.\\
Costs $1$ and $1+2^{-54}$ & Strictly different keys & Both compatibility costs round to 1.\\
Exact decimal midpoint 1.005 at 2 places & Nearest-even returns 1.00 & No binary64 conversion is part of this decimal quantization.\\
\bottomrule
\end{tabularx}
\caption{Selected conformance witnesses. Travel sequences can be embedded in a depot-to-depot route with zero remaining arcs and sufficiently wide other windows. These are designed examples, not measured BKS feasibility flips.}
\end{table}

Before a production implementation claims conformance, it \should{} compare with the reference on full retained-candidate corpora and generated boundary cases, then record its toolchain and independent-process reproduction. Cross-platform claims require actual cross-platform evidence. Bitwise agreement for the materialized inputs and rounded profiles, and exact agreement for the exact outputs, are stronger checks than matching formatted decimal displays. Metamorphic checks \should{} include route-list permutation invariance for exact profiles and scaled-integer/rational equivalence when eligibility is proved.

The execution record \must{} distinguish raw instance identity, materialized matrix or ATF identity, evaluator source identity and runtime identity. A successful rerun under a frozen environment supports reproducibility of that run. It does not establish unmeasured equivalence with another profile, another platform or an unavailable upstream evaluator.

\section{Adopting a new benchmark revision}\label{sec:migration}

A family adopting an evaluation profile publishes the specification and fixture revisions, the objective and units, and the manifests identifying its instances and their materialized inputs. A result remains associated with the family revision, evaluator and profile that produced it. A reader \may{} display an unversioned legacy record with its original attribution. Before admitting its routes to the new ranking, the store \must{} evaluate them under the new profile. A missing profile identifier \mustnot{} be interpreted as the current default.

A change in arithmetic can reverse the order of candidates, even when their displayed costs agree (Appendix~\ref{sec:evidence}). Preserve the old results and source records, evaluate every retained candidate under the new revision and rank feasible candidates by their complete objective key. If displaced candidates are missing, report the extent of the available history. For example, the SINTEF comparison identifies two official route sets that improve on the stored incumbents and a separate unresolved association between a published table and a route file. Preserve the downloaded files when correcting identifiers in derived records.

A migration report identifies the old and new revisions, coverage and missing inputs, changes in feasibility and ordering, differences at display precision, and the status of historical proof claims. A proof concerns the objective and feasibility rules under which it was established. Re-evaluating the routes supplies an incumbent value for the new revision. A claim of optimality under the new revision also requires a valid bound for that objective.

\section{Family coverage and result interchange}\label{sec:open}
% Internal decision record: /home/onyr/code/phd/tdvrptw-workspace/reports/plans/21-mamut-routing-lib-v013-plan.md, D13/D14 and OQ4. User ruling recorded on 2026-10-26 in session s121.

\paragraph{Family manifests.}
Each manifest identifies the files or regeneration inputs, the objective, units and fleet bounds. An integer or decimal-grid profile requires eligibility across the full instance. For example, a decimal value with 3 fractional digits admits an integer representation at scale 1000 under decimal interpretation. Reading that token as binary64 can yield a different rational number (Section~\ref{sec:profiles}). For large Mamut cases with unavailable road data, evaluation requires the missing sidecars or reproducibly regenerated values verified against the matrix digest.

For TD families, reproducibility depends on how the knots were obtained from stored ATFs, IGP speed profiles or a road graph and traffic overlay. To reproduce the knots, the family annex identifies the input revisions, units, horizon, preprocessing parameters and checksum encoding for each materializer. A hash of an ATF file that includes metadata cannot by itself serve as a digest of the knot sequence across different file encodings.

\paragraph{Result interchange.}
Section~\ref{sec:results} gives the proposed metadata fields for exact values and proof claims. A client that requires a numeric compatibility \id{cost} can consume an exact result only if its projection is finite. Family adoption must establish whether the accepted domain guarantees that bound or whether clients support the nullable projection. For \id{td-binary64-exact-v1}, the route-level format must associate each exact duration and optional departure witness with its canonical route. The Nairobi example in Section~\ref{sec:td-core} provides a round-trip case for the exact witness and its display.

\section*{Authorship and review}
% Internal editorial and decision provenance: /home/onyr/code/phd/phd_note_vault/docs/WORKLOG/2026-10-20.md, 2026-10-21.md, 2026-10-22.md and 2026-10-23.md. Session record: /home/onyr/code/phd/tdvrptw-workspace/reports/sessions/2026-10-23-session-121.md. Source identities: ../provenance.json.
Author order and ORCID identifiers follow the project citation record \autocite{pichonMAMUTRouting2026}. AI assisted with source inspection, writing and preparation of the executable companion. Coauthor review is pending.
